% ============================================================================
% บทที่ 6 — ตรวจสอบและแก้ปัญหาตารางสอบ
% ============================================================================
\mychapter{ตรวจสอบและแก้ปัญหาตารางสอบ}

\noindent บทนี้อธิบายการอ่านรายการข้อควรตรวจสอบ การสืบหาต้นตอของปัญหา และการแก้ไขข้อมูลต้นทาง

ระบบมีหน้าที่ชื่อคล้ายกันสามหน้า ใช้แยกกันดังนี้:

\begin{center}
\begin{tabular}{L{4cm}L{10cm}}
\toprule
\textbf{หน้า} & \textbf{ใช้เมื่อไร} \\
\midrule
ตรวจสอบข้อมูล & ดูผลหลังการนำเข้า ก่อนจัดตาราง \\
ข้อควรตรวจสอบ & ดูรายการปัญหาที่ตรวจพบ \\
ตรวจสอบ & ดูรายงานคุณภาพของตารางหลังจัด \\
\bottomrule
\end{tabular}
\end{center}

\section{ความรุนแรงและแหล่งที่มาของปัญหา}

ระบบแสดงปัญหาในหน้า \textbf{ข้อควรตรวจสอบ} แต่ละปัญหาระบุความรุนแรง
(\textbf{ข้อผิดพลาด} ซึ่งทำให้การจัดตารางไม่สมบูรณ์ หรือ \textbf{คำเตือน}
ซึ่งควรทราบแต่อาจไม่ขัดขวางการจัดตาราง) แหล่งที่มา
(\texttt{ต้นทาง} หมายถึงมาจากไฟล์ข้อมูลต้นทาง) และคำอธิบายสั้น ๆ
ใช้ช่องค้นหา \textbf{ค้นหาวิชาหรือปัญหา} และตัวกรองความรุนแรง/แหล่งที่มา
เพื่อหาปัญหาที่ต้องการ

\section{การขยายรายละเอียดปัญหา}

กดที่แถวของปัญหาเพื่อขยายรายละเอียด

\needspace{9\baselineskip}
\section{ตัวอย่าง: ปัญหาข้อมูลต้นทาง}

ตัวอย่างเช่น ปัญหา \textbf{ไม่พบวิชาในข้อมูลปัจจุบัน} สำหรับรหัส
\course{060999999} จะแสดงข้อความ:

\begin{center}
\textbf{เงื่อนไขอ้างถึงวิชาที่ไม่มีในไฟล์ที่จะจัดอัตโนมัติ}
\end{center}

พร้อมหลักฐานจากไฟล์ต้นทาง \file{rules.csv} \texttt{· Sheet1 · แถว 3 · คอลัมน์ 3}

\begin{expected}
เมื่อขยายรายละเอียด จะเห็นว่าปัญหามาจากเงื่อนไขที่อ้างถึงรหัสวิชาที่ไม่มีอยู่ใน
ไฟล์ที่จะจัดอัตโนมัติ วิธีแก้คือแก้ไขไฟล์ต้นทางแล้วนำเข้าใหม่อีกครั้ง
\end{expected}

\manualfigure{docs/usage-report/screenshots/figures/F20.png}{0.95\textwidth}{%
  ปัญหา \textbf{ไม่พบวิชาในข้อมูลปัจจุบัน} สำหรับรหัส \course{060999999}
  เมื่อขยายรายละเอียด พร้อมหลักฐานตำแหน่งในไฟล์ต้นทาง}

\manualfigure{docs/usage-report/screenshots/figures/F22.png}{0.9\textwidth}{%
  ปัญหา \textbf{แถวข้อมูลวิชาไม่ถูกต้อง} ในข้อมูลต้นทาง เมื่อขยายรายละเอียด
  พร้อมหลักฐานตำแหน่งในไฟล์ต้นทาง}

\section{การแก้ไขปัญหาข้อมูลต้นทาง}

เมื่อพบปัญหาที่เกิดจากข้อมูลต้นทาง ให้ขยายรายละเอียดของปัญหาก่อน
รายละเอียดจะแสดงไฟล์ต้นทาง ชื่อชีต หมายเลขแถว และหมายเลขคอลัมน์
ของข้อมูลที่เป็นต้นตอ พร้อมปุ่ม \texttt{ดูวิชา} เพื่อเปิดรายละเอียดของวิชาที่เกี่ยวข้อง
หมายเลขแถวและคอลัมน์นับตามตำแหน่งจริงในไฟล์เช่นเดียวกับโปรแกรมตารางคำนวณ
คือนับแถวหัวตารางเป็นแถวที่ 1 และนับคอลัมน์จากซ้ายเป็นคอลัมน์ที่ 1
สำหรับไฟล์ CSV ชื่อชีตจะแสดงเป็น \texttt{Sheet1}
จากนั้น:

\begin{steps}
  \item เปิดไฟล์ต้นทางในโปรแกรมตารางคำนวณ
  \item แก้ไขข้อมูลตามที่ระบุ (เช่น เพิ่มรหัสวิชาที่หาย หรือแก้กลุ่มนักศึกษา)
  \item บันทึกไฟล์ต้นทาง
  \item กลับมาที่ระบบแล้วนำเข้าไฟล์ใหม่อีกครั้ง
\end{steps}

หลังการนำเข้าใหม่ ปัญหาที่แก้ไขแล้วจะหายไปจากรายการข้อควรตรวจสอบ

\section{ตัวอย่าง: ปัญหากลุ่มนักศึกษาขาดหาย}

ปัญหากลุ่มนักศึกษาขาดหายจะทำให้ระบบไม่สามารถตรวจสอบการทับซ้อนของนักศึกษาได้
และอาจทำให้วิชานั้นไม่ถูกจัดตาราง ตัวอย่างสถานการณ์ก่อนการแก้ไข:

\manualfigure{docs/usage-report/screenshots/figures/F21.png}{0.95\textwidth}{%
  ตัวอย่างก่อนแก้ไข: หน้าภาพรวมแสดงวิชาที่ยังไม่ได้จัด เนื่องจากปัญหาข้อมูลต้นทาง
  (กลุ่มนักศึกษาขาดหาย) โดยกลางภาคและปลายภาคจัดได้เพียง 10/11 วิชา}

\begin{steps}
  \item ขยายรายละเอียดของปัญหาเพื่อดูวิชาที่ได้รับผลกระทบ
  \item เปิดไฟล์ต้นทางและตรวจสอบคอลัมน์ \texttt{studentGroups}
  \item เพิ่มกลุ่มนักศึกษาที่ขาดหาย
  \item นำเข้าไฟล์ใหม่อีกครั้ง
\end{steps}

\begin{expected}
หลังการแก้ไข วิชาที่ยังไม่ได้จัดจะถูกจัดตารางได้
จำนวนบนหน้าภาพรวมเปลี่ยนจาก 10/11 เป็น 11/11 และข้อผิดพลาดจะหายไป
\end{expected}

\section{การอ่านหน้ารายงานการตรวจสอบ}

หน้ารายงานการตรวจสอบ (\texttt{ตรวจสอบ}) แสดงข้อมูลคุณภาพของตารางสอบ
เช่น จำนวนวิชาที่จัดแล้ว จำนวนวันที่มีการสอบ จำนวนวิชาที่สอบพร้อมกันสูงสุด
และผลการตรวจสอบความถูกต้อง สำหรับข้อมูลที่ถูกต้องครบถ้วน
หน้ารายงานจะแสดง \textbf{0 ข้อผิดพลาด} และ \textbf{0 คำเตือน}
สามารถดาวน์โหลดรายงานตรวจสอบฉบับเต็มเพื่อส่งให้ผู้ดูแลระบบได้ (ดูบทที่ 7)

\begin{note}
การที่รายการปัญหาแสดงจำนวน 0 ไม่ได้หมายความว่าตารางสอบสมบูรณ์แบบเสมอไป
ควรตรวจสอบจำนวนวิชาที่ยังไม่ได้จัดและความถูกต้องของข้อมูลด้วย
\end{note}

\manualfigure{docs/usage-report/screenshots/crops/F23b.png}{0.9\textwidth}{%
  รายการตรวจสอบในหน้ารายงานการตรวจสอบ}

\seealso{บทที่ 5 สำหรับการแก้ไขเวลาสอบ และบทที่ 7 สำหรับการดาวน์โหลดผลลัพธ์}
